\documentclass[12pt]{article}

% Packages
\usepackage[
  paperwidth=8.5in,
  paperheight=11in,
  margin=0.5in
]{geometry}
\usepackage{amsmath, amssymb}
\usepackage{booktabs}
\usepackage{tabularx}
\usepackage{graphicx}
\usepackage{hyperref}


\newcommand{\tacred}{FS-TACRED}
\newcommand{\qwensmall}{Qwen3-4B}
\newcommand{\gemmasmall}{Gemma3-4B}

\begin{document}

\paragraph{Idea of a new approach for prompt optimization}
\begin{itemize}
    \item Optimize a {set of diverse prompts} instead of a single prompt
    \item Organize prompts into {clusters} capturing different aspects of the task, by splitting errors/feedback (with some overlap) across clusters.
    \item Perform {local optimization} within clusters (error analysis, mutation, crossover)
    \item Allow inter-cluster crossover at the end.
    \item Moreover, use {global error analysis} to guide overall improvement
    \item Combine outputs from multiple final prompts to make decision. Here are some ideas to optimize the inference cost:
    \begin{itemize}
        \item Use only a {subset of the prompts} to reduce inference cost.
        \item Start with the {best prompt}, and iteratively add others until we are confident in the prediction.
        \item Use a {router or classifier} to map inputs to the most suitable prompt(s).
    \end{itemize}
\end{itemize}


\paragraph{Other updates}
I evaluated several optimized prompts (from earlier experiments) on the FS-TACRED test set, 
and they show improvement over the baseline prompt. 
However, prompts optimized without feedback show a drop in performance,
because they were optimized directly on the validation set (overfit), whereas other approaches obtain feedback from the training set.

% \bibliographystyle{plain}  
% \bibliography{custom}    

\end{document}
